\chapter{Duplicar una sucesión convergente}

\section{El desafío de escalar en la fábrica}

Tras conquistar la sucesión constante, subamos el nivel: si una sucesión
converge, entonces duplicar esa sucesión también converge, y converge al doble
del límite original.

Imaginemos que el Maquinista recibe un encargo del Ingeniero: «Por favor,
duplica todas las longitudes, pero mantén los mismos estándares de calidad».
¿Cómo debería responder el Maquinista?

Si el Ingeniero exige que las longitudes duplicadas estén dentro de $\varepsilon$
de $2L$, el Maquinista no puede simplemente pedir que el proceso original
cumpla la tolerancia original $\varepsilon$, porque
\[
  |2a(n) - 2L| < 2\varepsilon,
\]
no $\varepsilon$. En cambio, debemos ser más ingeniosos. ¿Se te ocurre qué hacer?

Exacto: el $\varepsilon$ del proceso original es arbitrario, ¡así que podemos
jugar con él! Si logramos que el proceso original garantice una salida dentro
de $\varepsilon/2$ de $L$, entonces al duplicar esa salida estaremos dentro de
$\varepsilon$ de $2L$.

\medskip
\noindent\textbf{Idea clave:} al escalar por una constante, debemos escalar
las exigencias de tolerancia de forma inversa.

\section{La configuración matemática}

Supongamos que tenemos:
\begin{itemize}
  \item Una sucesión $a : \mathbb{N} \to \mathbb{R}$ que converge a $L$.
  \item Una sucesión $b : \mathbb{N} \to \mathbb{R}$ con la propiedad de que
        $b(n) = 2a(n)$ para todo $n$.
\end{itemize}
Queremos demostrar que $b$ converge a $2L$.

\subsection{Idea Clave: Escalado Inverso de la Tolerancia}

La observación crucial es que:
\[
  |b(n) - 2L| = |2a(n) - 2L| = 2|a(n) - L|.
\]
Así, si queremos $|b(n) - 2L| < \varepsilon$, necesitamos
$2|a(n) - L| < \varepsilon$, es decir:
\[
  |a(n) - L| < \frac{\varepsilon}{2}.
\]
¡Esto es exactamente lo que nos proporciona la convergencia de $a$!

\section{Nuevas herramientas que necesitarás}

\subsection{Valor absoluto de productos}

Necesitarás el nuevo teorema \texttt{abs\_mul}, que establece que para
cualesquiera números reales $x$ e $y$:
\[
  |xy| = |x||y|.
\]
Para usar este teorema puede resultar conveniente crear una nueva hipótesis
con \texttt{have} y luego reescribir con esa hipótesis. Es decir, puedes
escribir:

\begin{minted}{lean}
have NuevoHecho : | Algo * OtraCosa | =
    | Algo | * | OtraCosa | :=
  by apply abs_mul
\end{minted}

y después \mintinline{lean}{rewrite [NuevoHecho]} sustituirá
\mintinline{lean}{|Algo * OtraCosa|} por
\mintinline{lean}{|Algo| * |OtraCosa|} (ya sea en el Objetivo o en una
hipótesis, si así se especifica).

\section{Enfoque estratégico}

\begin{enumerate}
  \item Despliega la definición de convergencia para el objetivo.
  \item Dado cualquier $\varepsilon > 0$, usa la convergencia de $a$ con
        tolerancia $\varepsilon/2$. Puede ser útil demostrar por separado la
        desigualdad $0 < \varepsilon/2$: ¿qué táctica crees que ayudaría aquí?
  \item Extrae el testigo $N$ de la convergencia de $a$.
  \item Usa el mismo $N$ para tu sucesión $b$.
  \item Aplica la hipótesis de escalado y usa \texttt{abs\_mul} para relacionar
        $|b(n) - 2L|$ con $|a(n) - L|$.
  \item Usa la cota de convergencia de $a$ para concluir.
\end{enumerate}

\section{Solución en Lean}

\begin{minted}{lean}
-- Demostración en lenguaje natural
-- ================================

-- Sea bₙ la suceción definida por
--    bₙ = 2aₙ                                                        (1)
-- Tenemos que demostrar que para cada ε > 0, existe un N ∈ ℕ tal que
--    ∀ n ≥ N, |bₙ - 2L| < ε                                          (2)
-- Puesto que aₙ converge a L, existe un N ∈ ℕ tal que
--    ∀ n ≥ N, |aₙ - L| < ε/2                                         (3)
-- Veamos que N también cumple (2). En efecto, sea n ≥ N. Entonces
--    |bₙ - 2L| = |2aₙ - 2L|     [por (1)]
--              = 2|aₙ - L|
--              < 2ε/2           [por (3)]
--              = ε

-- Demostraciones con Lean 4
-- =========================

import Mathlib.Data.Real.Basic
import Mathlib.Tactic
import Analisis_en_Lean.C2_El_calculo_de_pi

variable {a b : ℕ → ℝ}
variable {L : ℝ}

-- 1ª demostración
-- ===============

example
  (ha : LimSuc a L)
  (hb : ∀ n, b n = 2 * a n)
  : LimSuc b (2 * L) :=
by
  intro ε hε
  -- ε : ℝ
  -- hε : ε > 0
  -- ⊢ ∃ N, ∀ n ≥ N, |b n - 2 * L| < ε
  obtain  ⟨N, hN⟩ := ha (ε / 2) (by grind)
  -- N : ℕ
  -- hN : ∀ n ≥ N, |a n - L| < ε / 2
  use N
  -- ⊢ ∀ n ≥ N, |b n - 2 * L| < ε
  intro n hn
  -- n : ℕ
  -- hn : n ≥ N
  -- ⊢ |b n - 2 * L| < ε
  grind

-- 2ª demostración
-- ===============

example
  (ha : LimSuc a L)
  (hb : ∀ n, b n = 2 * a n)
  : LimSuc b (2 * L) :=
by
  intro ε hε
  -- ε : ℝ
  -- hε : ε > 0
  -- ⊢ ∃ N, ∀ n ≥ N, |b n - 2 * L| < ε
  obtain  ⟨N, hN⟩ := ha (ε / 2) (by grind)
  -- N : ℕ
  -- hN : ∀ n ≥ N, |a n - L| < ε / 2
  use N
  -- ⊢ ∀ n ≥ N, |b n - 2 * L| < ε
  intro n hn
  -- n : ℕ
  -- hn : n ≥ N
  -- ⊢ |b n - 2 * L| < ε
  calc |b n - 2 * L|
       = |2 * a n - 2 * L| := by grind
     _ = |2 * (a n - L)|   := by grind
     _ = 2 * |a n - L|     := by grind
     _ < 2 * ε / 2         := by grind
     _ = ε                 := by grind

-- 3ª demostración
-- ===============

example
  (ha : LimSuc a L)
  (hb : ∀ n, b n = 2 * a n)
  : LimSuc b (2 * L) :=
by
  intro ε hε
  -- ε : ℝ
  -- hε : ε > 0
  -- ⊢ ∃ N, ∀ n ≥ N, |b n - 2 * L| < ε
  obtain  ⟨N, hN⟩ := ha (ε / 2) (half_pos hε)
  -- N : ℕ
  -- hN : ∀ n ≥ N, |a n - L| < ε / 2
  use N
  -- ⊢ ∀ n ≥ N, |b n - 2 * L| < ε
  intro n hn
  -- n : ℕ
  -- hn : n ≥ N
  -- ⊢ |b n - 2 * L| < ε
  calc |b n - 2 * L|
       = |2 * a n - 2 * L| := by rw [hb n]
     _ = |2 * (a n - L)|   := by simp only [mul_sub]
     _ = |2| * |a n - L|   := by simp only [abs_mul]
     _ = 2 * |a n - L|     := by simp only [abs_two]
     _ < 2 * (ε / 2)       := by simp only [ mul_lt_mul_of_pos_left, hN n hn, two_pos]
     _ = ε                 := by rw [mul_div_cancel₀ ε two_ne_zero]

-- 4ª demostración
-- ===============

example
  (ha : LimSuc a L)
  (hb : ∀ n, b n = 2 * a n)
  : LimSuc b (2 * L) :=
by
  intro ε hε
  -- ε : ℝ
  -- hε : ε > 0
  -- ⊢ ∃ N, ∀ n ≥ N, |b n - 2 * L| < ε
  obtain  ⟨N, hN⟩ := ha (ε / 2) (half_pos hε)
  -- N : ℕ
  -- hN : ∀ n ≥ N, |a n - L| < ε / 2
  use N
  -- ⊢ ∀ n ≥ N, |b n - 2 * L| < ε
  intro n hn
  -- n : ℕ
  -- hn : n ≥ N
  -- ⊢ |b n - 2 * L| < ε
  calc |b n - 2 * L|
       = |2 * a n - 2 * L| := congr_arg (|· - 2 * L|) (hb n)
     _ = |2 * (a n - L)|   := congr_arg abs (mul_sub 2 (a n) L).symm
     _ = |2| * |a n - L|   := abs_mul 2 (a n - L)
     _ = 2 * |a n - L|     := congr_arg (· * |a n - L|) abs_two
     _ < 2 * (ε / 2)       := mul_lt_mul_of_pos_left (hN n hn) two_pos
     _ = ε                 := mul_div_cancel₀ ε two_ne_zero

-- Lemas usados
-- ============

variable (x y z : ℝ)

#check (abs_mul x y : |x * y| = |x| * |y|)
#check (abs_two : |(2 : ℝ)| = 2)
#check (half_pos : 0 < x → 0 < x / 2)
#check (mul_div_cancel₀ x : y ≠ 0 → y * (x / y) = x)
#check (mul_lt_mul_of_pos_left : x < y → 0 < z → z * x < z * y)
#check (mul_sub x y z : x * (y - z) = x * y - x * z)
#check (two_ne_zero : 2 ≠ 0)
#check (zero_lt_two : 0 < 2)

-- 3ª demostración
-- ===============

example
  (ha : LimSuc a L)
  (hb : ∀ n, b n = 2 * a n)
  : LimSuc b (2 * L) :=
by
  change ∀ ε > 0, ∃ N : ℕ, ∀ n ≥ N, |b n - 2 * L| < ε
  -- ⊢ ε > 0 → ∃ N, ∀ n ≥ N, |b n - 2 * L| < ε
  intro ε hε
  -- ε : ℝ
  -- hε : ε > 0
  -- ⊢ ∃ N, ∀ n ≥ N, |b n - 2 * L| < ε
  change ∀ ε₁ > 0, ∃ N₁ : ℕ, ∀ n ≥ N₁, |a n - L| < ε₁ at ha
  -- ha : ∀ ε₁ > 0, ∃ N₁, ∀ n ≥ N₁, |a n - L| < ε₁
  specialize ha (ε / 2)
  -- ha : ε / 2 > 0 → ∃ N₁, ∀ n ≥ N₁, |a n - L| < ε / 2
  have hε2 : 0 < ε / 2 := by bound
  specialize ha hε2
  -- ha : ∃ N₁, ∀ n ≥ N₁, |a n - L| < ε / 2
  choose N hN using ha
  -- N : ℕ
  -- hN : ∀ n ≥ N, |a n - L| < ε / 2
  use N
  -- ⊢ ∀ n ≥ N, |b n - 2 * L| < ε
  intro n hn
  -- n : ℕ
  -- hn : n ≥ N
  specialize hb n
  -- hb : b n = 2 * a n
  rewrite [hb]
  -- ⊢ |2 * a n - 2 * L| < ε
  clear hb hε2
  have h1 : 2 * a n - 2 * L = 2 * (a n - L) := by ring_nf
  rewrite [h1]
  -- ⊢ |2 * (a n - L)| < ε
  clear h1
  have h2 : |2 * (a n - L)| = |2| * |a n - L| := by apply abs_mul
  rewrite [h2]
  -- ⊢ |2| * |a n - L| < ε
  clear h2
  specialize hN n hn
  -- hN : |a n - L| < ε / 2
  clear hn
  norm_num
  -- ⊢ 2 * |a n - L| < ε
  bound
\end{minted}

\section{Prueba en lenguaje natural}

Demos un paso atrás respecto a la prueba formal en Lean y comprendamos
en castellano lo que acabamos de demostrar.

\subsection{Enunciado del teorema}

\begin{theorem}
Si una sucesión de números reales converge a algún límite, entonces la sucesión
formada al duplicar cada término converge al doble del límite original.
\end{theorem}

\subsection{Prueba}

\begin{proof}
Supongamos que la sucesión $a_n$ converge a $L$, y queremos demostrar que
$b_n = 2a_n$ converge a $2L$.

Por definición, debemos demostrar que para cualquier tolerancia $\varepsilon > 0$,
podemos encontrar un punto $N$ tal que para todo $n \geq N$ se tenga
$|b_n - 2L| < \varepsilon$.

La idea clave es la siguiente: como $a_n$ converge a $L$, podemos hacer
$|a_n - L|$ arbitrariamente pequeño. Concretamente, podemos encontrar un $N$
tal que $|a_n - L| < \varepsilon/2$ para todo $n \geq N$.

Ahora bien, para cualquier $n \geq N$:
\begin{align}
  |b_n - 2L| &= |2a_n - 2L| \\
             &= |2(a_n - L)| \\
             &= 2|a_n - L| \\
             &< 2\frac{\varepsilon}{2} \\
             &= \varepsilon
\end{align}
Por tanto, $b_n$ converge a $2L$, completando la prueba.
\end{proof}

\section{Por qué importa este patrón}

Esta prueba introduce la importante técnica del \emph{escalado inverso de
tolerancias}. Cuando se escala una sucesión por una constante $c$, es necesario
escalar las exigencias de tolerancia por $1/c$. Este principio reaparecerá
cuando estudies:

\begin{itemize}
  \item Productos de sucesiones (donde ambos factores contribuyen al error).
  \item Derivadas (donde la definición de límite involucra escalado por $h$).
  \item Integración (donde las sumas de Riemann involucran escalado por los
        anchos de partición).
\end{itemize}

¡La capacidad de gestionar cómo las constantes afectan a la convergencia es
fundamental en todo el análisis!

\subsection{La estrategia elegante}

La prueba utilizó el \emph{principio de escalado inverso}:

\begin{enumerate}
  \item \textbf{Inversión de la tolerancia:} para lograr tolerancia $\varepsilon$
        en la salida duplicada, se exige tolerancia $\varepsilon/2$ en la
        entrada original.
  \item \textbf{Factorización algebraica:} $2a_n - 2L =
        2(a_n - L)$ reveló la estructura.
  \item \textbf{Escalado del valor absoluto:} $|2x| = |2||x|$,
        seguido de \texttt{norm\_num}, convirtió la forma factorizada en la
        cota necesaria.
  \item \textbf{Aritmética lineal:} la cota final combinó
        $2|a_n - L| < 2(\varepsilon/2) = \varepsilon$.
\end{enumerate}

%%% Local Variables:
%%% mode: LaTeX
%%% TeX-master: "Analisis_en_Lean"
%%% End:
